- montrer que bien que les frameworks partagent les même concepts, ils ne suivent  pas la même logiques de languages et d'architecture. donc développer dans un framework demande de le comprendre, de comprendre son archi de conception pour pouvoir véritablement l'utiliser. Exemple avec la mémoire, l'orchestration, les types d'agent chez adk etc.

- peut-être prendre aussi en compte le temps de prise en main ? 

- on va déjà pendre en compte : 
    - temps de réalisation de la tache
    - Est ce que la tache a été accompli ou pas ? 
    - Durée

- peut-être avoir le public de novice à qui on va donner un essaie pour la prise en main du truc. et comparer aussi : 
    - le temps de prise en main -> le temps pour réaliser le truc de base en se basant sur la documentation.



Voici ma proposition de RQs : 
 \begin{itemize}
     \item \textbf{RQ1 :} Est ce que les FM (de l'agent, du module et du multiagent) expriment les concepts necessaires pour décrire des workflows multiagents  et empêchent les incohérences ?
     \item \textbf{RQ2 :} Quelle est la proportion des concepts Gear traduite vers chaque framework (sans perte, avec perte, approximations et exact) et quel est l'impact de ces écarts sur les artefacts générés
     \item \textbf{RQ3 :} Par rapport au développement natif (CrewAI/ADK), Gear réduit‑il le temps, la quantité de code/artefacts, les tokens, et comment ces gains évoluent‑ils quand la taille/complexité du workflow augmente pour différents profils (débutants, industriels, chercheurs) ? \textit{je crois que ce serait bien de regarder à ce niveau le time‑to‑run, taux d’erreurs, stabilité}
     \item \textbf{RQ4 :} Dans quelle mesure Gear permet‑il une comparaison systématique et équitable des frameworks (CrewAI vs ADK) à partir d’un même modèle Gear
 \end{itemize}

 je crois qu'on doit aussi ajouter deux choses dans la partie Threats to validity :
 \begin{itemize}
     \item effort pour ajouter un nouveau connector et résister aux évolutions des frameworks
     \item capacité à relier une erreur/écart dans le code généré à une feature UVL
 \end{itemize}

 